\clearpage\thispagestyle{empty}
\pdfbookmark[0]{Quellen zur Besprechung von Seyffers Schrift über die Länge von München}{seyffer1811-sources}
\begin{center}\textbf{Quellen zur Besprechung von Seyffers Schrift über die Länge von München}\end{center}
\noindent Diese KI-übersetzte Fassung in modernem Deutsch folgt Carl Friedrich Gauß, \textit{Werke}, Band~VI (Göttingen, 1874), S.~533--534. Die Besprechung erschien zuerst unter München in den \textit{Göttingischen gelehrten Anzeigen}, Stück~78, 18. Mai 1811, S.~773--775.
\par\smallskip\noindent Der lateinische Titel lautet in der Werkausgabe \textit{Super longitudine geographica speculae astronomicae regiae, quae Monachii est, ex triginta septem defectionibus Solis observatis et ad calculos revocatis nunc primum definita}. Seyffers Titelblatt von 1810 bezeichnet die Schrift zusätzlich als \textit{Commentatio altera}. Das digitalisierte Exemplar stammt aus der Bayerischen Staatsbibliothek.
\par\smallskip\noindent Der Erstdruck verwendet bei den beiden Längenangaben einfache und doppelte Striche. In der Werkausgabe stehen hochgestellte m und s sowie die Zeichensetzung \textit{Brünn. Berlin}. Die Übertragung folgt der Werkausgabe.
\par\smallskip\noindent \href{https://gdz.sub.uni-goettingen.de/id/PPN236007467}{Text der Werkausgabe}. Erstdruck, \href{https://images.p.sub.uni-goettingen.de/iiif/image/gdz:PPN319721507_1811:00000805/full/1650,/0/default.jpg}{S.~773}, \href{https://images.p.sub.uni-goettingen.de/iiif/image/gdz:PPN319721507_1811:00000806/full/1650,/0/default.jpg}{S.~774} und \href{https://images.p.sub.uni-goettingen.de/iiif/image/gdz:PPN319721507_1811:00000807/full/1650,/0/default.jpg}{S.~775}. \href{https://api.digitale-sammlungen.de/iiif/image/v2/bsb10723478_00009/full/1550,/0/default.jpg}{Seyffers Titelblatt}.
\par\smallskip\noindent Übersetzung und Satz mit KI erstellt.
\clearpage
